\section{Discussion}
\label{sec:discussion}

\subsection{The applicability window is a result}

The gate of Eq.~\ref{eq:gate} and the multiplier of Eq.~\ref{eq:gain} together
form a deployment rule that can be evaluated before building anything: a second
KV tier pays when the latency target exceeds the full-pool step time, and it
pays in proportion to $\HBM/\Ckv$ --- how much of device memory the weights
consumed. Both exclude configurations the field routinely evaluates on. An 8B
model on a 96\,GB device leaves so much memory after weights that there is
nothing to buy, and a latency target below the device's own step time makes the
scheduler latency-bound rather than capacity-bound, where offloading is
strictly harmful.

We state the window as a finding rather than a caveat because the alternative
is worse: a technique whose applicability is unstated gets evaluated where it
happens to win.

\subsection{Generality across memory tiers and attention backends}
\label{sec:discussion:general}

Nothing in the design is specific to HBM-plus-DRAM or to one attention kernel.
The system needs three things from a substrate: a memory that is larger and
slower than the device's, a path to it whose bandwidth can be measured, and an
attention implementation that can read a block table. Everything else --- how
many requests to place on the slower tier, which ones, when to move them --- is
derived at runtime from quantities the system measures itself.

\paragraph{More than two tiers.} The placement rule ranks requests by a scalar
affinity and the admission controller enforces a per-step byte budget on one
link. A third tier (CXL-attached memory, or NVMe) adds a second budget and a
second affinity threshold; the structure is unchanged, and the parameters that
distinguish tiers --- achieved bandwidth, and the per-request bytes a step must
move --- are already measured rather than configured. What does not survive
unchanged is the hiding argument: a tier whose transfer cannot fit inside a
layer's compute window must be treated as a capacity tier for idle requests
rather than a bandwidth tier for active ones, and Eq.~\ref{eq:gate} is what
decides which it is.

\paragraph{Different attention capabilities.} Our host tier stages KV through a
rolling device buffer. A zero-copy kernel that reads host memory directly
removes that buffer, and our backend interface exists so it can be substituted
without touching the policy. The two are not ordered by quality but by
operating point (\S\ref{sec:discussion:coherence}): staging buys a layer-sized
hiding window at the cost of two extra device-memory passes, and which wins
depends on whether the per-step transfer fits inside a layer's compute. Because
the offload ratio here is chosen by admission rather than fixed by the design,
that choice is ours to make from a measurement rather than a constraint we
inherit.

\paragraph{What must be recalibrated, and what must not.} Moving to a new
platform requires re-measuring four quantities: achieved interconnect
bandwidth, achieved device-memory bandwidth for weights and for paged KV, and
the fixed per-step cost of the host path. The closed-form rules do not change;
their inputs do. This is why we report the fixed cost separately from the
per-byte cost rather than as a single slope --- only one of the two shrinks
when the interconnect improves, and a system calibrated on their sum will
mispredict its own behaviour on the next generation of hardware.

\subsection{Anti-scaling}

The gain shrinks as device memory grows, and for our primary model it does not
merely shrink: Qwen2.5-32B on a 144\,GB device leaves $\Ckv \approx 73$\,GB, at
which point the workload is no longer capacity-bound and the applicability gate
stops admitting entirely. This is the same structural decay we identify in
compute-offload systems, whose advantage fades as GPUs strengthen, and it would
be dishonest to cite that against them without stating it about ourselves.

Two things bear on the trend. Weights are growing faster than device memory
--- mixture-of-experts and 405B-class models push $\HBM/\Ckv$ up, not down ---
so the configurations where the technique pays are becoming more common rather
than less. And the gate is the deployment rule: a system that switches itself
off where it cannot help is behaving correctly, not failing.

\subsection{Composition with orthogonal techniques}

Methods that \emph{shrink} the KV cache compose with a system that
\emph{relocates} it. Quantizing KV scales the device-resident cohort and the
affordable host cohort together, leaving their ratio --- and so the value of a
second tier --- unchanged. Quantizing \emph{weights} does the opposite: it
returns memory to $\Ckv$ and removes the amplifier the technique depends on,
which is the same mechanism by which a larger device reduces the gain.
Sparse and low-rank methods reduce what must be stored at some cost in
exactness; a tier moves bytes without altering them, so the two address
different axes and could be combined.

\subsection{Limitations}

The prefill transport is emulated: KV reaches the host pool through a local
handover whose cost we measure (\S\ref{sec:impl:stash}), not through a
production prefill node's network path, which is deployment-dependent.
Speculative decoding is unsupported on the host path, which assumes one query
token per step. The pinned host mirrors are real memory, sized to the
configured maximum concurrency and context. We evaluate a single device;
tensor parallelism would introduce collectives into precisely the metric under
measurement. And the fixed per-step cost of the streaming path is a floor of
the current design (\S\ref{sec:impl:cost}) rather than a tuned result --- a
cross-layer batching of the transfer issue would reduce it, and we have not
built it.

\subsection{Staging versus zero-copy: two hiding windows}
\label{sec:discussion:coherence}

Our host tier stages: a layer's KV is copied from a pinned host mirror into a
device rolling slot before attention reads it. DirectKV's kernel instead reads
CPU-resident KV directly, with no staging buffer at all. It is tempting to read
our buffer as a tax that their kernel would simply remove. That reading is
wrong, and the reason is worth setting out because it determines which
mechanism belongs at which operating point.

\paragraph{The two designs hide the transfer behind different things.} A
zero-copy kernel overlaps its fetch with the arithmetic of the \emph{attention
tile it is already computing} --- warp-level pipelining, a window of one tile.
Decode attention's arithmetic intensity is the GQA group size, 4--8 FLOP/byte
against a hardware balance point near 116, so there is almost no arithmetic
there to hide behind: the kernel is bandwidth-bound on the fetch it is trying
to hide. Staging overlaps the fetch for layer $\ell{+}1$ with the \emph{entire
compute of layer $\ell$}, which includes the QKV, output and MLP projections.
Those are weight-bound and long --- about 21 of a 34\,ms step for Qwen-32B on
GH200. \textbf{The staging buffer is not overhead; it is what buys a
layer-sized hiding window instead of a tile-sized one.}

\paragraph{What staging costs, and hence where the crossover is.} Staging pays
two extra HBM passes per byte --- the write into the slot and the read out of
it --- plus the slot's own device footprint. Zero-copy pays neither. So the
comparison is not one design dominating the other but a crossover in the
\emph{offload ratio}:

\begin{itemize}
\item When the per-step transfer \textbf{fits inside the layer compute window},
      staging hides it almost entirely and pays only the extra HBM traffic,
      which at $\Xlink/\BW \approx 0.15$ is a small fraction of a transfer that
      is itself hidden. Staging wins.
\item When the transfer \textbf{exceeds that window}, it is exposed under either
      design, and zero-copy wins by avoiding the extra HBM passes and the
      buffer's footprint.
\end{itemize}

That is precisely why the two systems make opposite choices, and both are
right. DirectKV places 100\% of KV on the host for every request, so its
per-step transfer is the entire KV working set and can never fit a compute
window --- zero-copy is the correct mechanism for that design point, and its
CPU-memory-aware tiling is aimed at shifting the residual bottleneck back to
HBM. We place a \emph{selected subset} of requests on the host, and admission
sizes that subset so its transfer fits the window by construction
(\S\ref{sec:design:admission}). \textbf{The offload ratio is not a fixed
property of the system here; it is the quantity our policy controls} --- which
means the choice of mechanism follows from the policy rather than constraining
it.

\paragraph{Where a coherent link helps regardless: the write path.} One cost of
ours is not about hiding at all. Each step projects new $K,V$ for every
host-tier request and we copy them device-to-host to keep the mirror
authoritative. If those blocks were mapped host memory the kernel could write
them in place, removing the copy and the bookkeeping that keeps two copies
consistent. The effect is small in time --- an asynchronous version of that
copy measured 1.6\,ms --- but it removes a class of state rather than an amount
of time, and that state is where migration bugs live. DirectKV's Alg.~3 already
does exactly this, appending newly projected $K,V$ to CPU memory from inside
the kernel.

\paragraph{What coherence does not buy.} It does not make the transfer free; it
makes it implicit. Coherence is a guarantee about \emph{visibility}, not a
cheaper transport --- the same bytes cross the link at link bandwidth, fetched
when the kernel touches them. Nor does anything need to poll for completion:
nothing on the host consumes the KV, since the host copy exists so that the
\emph{GPU} can read it back later. Were a consumer introduced, the construction
is a release-ordered flag with system-scope fencing, and the ordering rather
than the notification is the hard part.

\subsection{Placement by hardware, policy by scheduler}

A further consequence, which we flag as future work because we cannot test it.
With a single coherent copy, demoting a request to the host tier becomes a
policy decision rather than a data movement, and promotion could be lazy and
partial --- the hardware's page migration relocating hot blocks into device
memory while the memory system supplies hints rather than issuing transfers.
That is an appealing division of labour, and to our knowledge nobody has
evaluated it for KV cache.

We are wary of it for one specific reason. Hardware page migration operates at
4\,KB or 2\,MB granularity and migrates on access, which is structurally the
same reactive, fine-grained policy whose migration overhead dominates in the
token-granularity placement literature~\cite{calkv}. Our own claim is that
request-granularity movement avoids that thrash because whole requests move
once, at controlled moments. Handing placement back to a fault-driven mechanism
gives that property up, and whether the hardware's migration is cheap enough to
survive it is an empirical question we have not answered.
